\documentclass{ltjsbook}
\usepackage{docmute} 
\usepackage{../../myStyle} 

\begin{document}
\chapter{確率モデルとベイズ推測}
\label{H13_Bayesian}


\section{ベイズ推定の基礎}

データ生成モデルが確率の言葉で記述される，というのはすでに述べたところです。推測統計というのがそもそも不良設定問題，すなわち少なすぎるヒントから未知なるものを当てるという状況に置かれた学問であり，この「わからないこと」を確率で表現するからです。この確率についての考え方として，ベイズの公式というのがあるのでした。ベイズの定理は次のように書き表されます。

$$ P(\theta|D) = \frac{P(D|\theta)P(\theta)}{P(D)} $$

ここで$P(X)$はXについての確率，という表現です。$P(A|B)$は条件付き確率というもので，Bが与えられた時のAの確率，という意味です。

ベイズの定理の用語は次の通りです。まず右辺の分子に注目しましょう。$P(D|\theta)$とありますが，ここで$D$はデータを表しています。$\theta$はデータを生む確率のパラメータです。$P(D|\theta)$はパラメータ$\theta$のもとでのデータ$D$が得られる程度を表すもの，という意味になります。これは\keyterm{尤度}{ゆうど}{likelihood}と呼ばれるもので，パラメータの関数になっているのでした。たとえば正規分布からデータが生成されると考えるのであれば，手元のデータがパラメータ$\theta$から出てくる可能性がどれぐらいあるのか，を表現していると言えるわけです。この尤度関数は，データを生み出すメカニズムの表現そのものです。

右辺の分子の第二項，$P(\theta)$はパラメータの確率です。正規分布からデータが生成される例で言えば，尤度はあるパラメータの値からデータが得られる程度を表現しているのですが，そのパラメータが出てくる確率はそもそもどの程度であるか，を考えているのです。これは別名\keyterm{事前分布}{じぜんぶんぷ}{prior distribution}とも呼ばれます。実際にデータが出てくる前の段階で，そもそもそのパラメータがどれほど出やすいかという確率を表現しており，これは今回のデータを取る前までの事前のデータや経験に基づいている，ということができます。

右辺の分母，$P(D)$はデータ全体が得られる確率であり，\keyterm{周辺尤度}{しゅうへんゆうど}{marginal likelihood}とか\keyterm{エビデンス}{えびでんす}{evidence}と呼ばれます。\keyterm{正規化定数}{せいきかていすう}{normalized constant}ということもあります。これについては理解しにくいところもあるかもしれませんが，後に述べる理由によって，ひとまず深く考える必要はありません。

左辺はこの計算の結果得られる\keyterm{事後分布}{じごぶんぷ}{posterior distribution}を表しています。$P(\theta|D)$はデータ$D$で条件づけられたパラメータ$\theta$の確率です。我々は確率モデルを作るわけですが，そのパラメータがどうなっているかは事前にはわかりません。が，データを取ることで「データが与えられたのであれば，パラメータはこうだよ」というのがわかるわけです。あるいは事前にパラメータはこのあたりにあるのではないか，と経験上考えていたとしても，それがデータを取ることによって更新される(確信が強くなったり，違うかもしれないと思い直したり)，ということを意味しています。

この式を言葉で表現し直すと，次のようになります。

$$ \mbox{事後分布} = \frac{\mbox{尤度}\times \mbox{事前分布}}{\mbox{周辺尤度}} = \mbox{尤度} \times \frac{事前分布}{周辺尤度}$$

\keytermL{尤度}{ゆうど}がメカニズムそのものだ，ということはすでに書きました。たとえば\keytermL{回帰分析}{かいきぶんせき}においても，尤度を計算できます。普通の回帰分析は，誤差が確率的に生じると考え，誤差以外のところは線形モデルで表現します。すなわち，従属変数$Y_i$に対して，独立変数$X_i$があったとすると，

\begin{equation*}
    \begin{array}{lll}
        Y_i & =    & \hat{Y}_i + e_i = \beta_0 + \beta_1 X_i + e_i \\
        e_i & \sim & N(0,\sigma)                                   \\
    \end{array}
\end{equation*}
より，
$$ Y_i \sim N(\beta_0 + \beta_1 X_i, \sigma) $$
となるのでした。ここで$N(\mu,\sigma)$は平均$\mu$，標準偏差$\sigma$の正規分布を表し，データ$Y_i$が正規分布から生成されているというモデルになっています。回帰分析の場合は\keytermL{最尤法}{さいゆうほう}で未知数$\beta_0,\beta_1,\sigma$を求めたりしましたが\footnote{あるいは\keyterm{最小二乗法}{さいしょうにじょうほう}{Least Square method}で求めるのですが，これは幸い最尤法の結果と一致します。正規分布を使った線形回帰モデルの場合，データの記述統計的値と確率モデルによる推定値が一致するので，気づかないまま推測統計学に足を踏み入れてしまえるのでした}，その名の通りこのモデルが示す尤度を最大にする未知数の求め方を最尤法というのでした。

ですが，尤度は確率を表すものではありません。確率とは非負の実数で，すべての確率を足し合わせる(あるいは積分する)と$1.0$にならなければなりませんが，尤度は足し合わせても1.0にならない数字なのです。ですから，「手元のデータがパラメータからでてくる\textbf{可能性}」という変な言い回しをしていました。「出てくる\textbf{確率}」とは言えないのですね。そこで，尤度を確率に置き換えたい，と考えたときに便利なのがベイズの定理なのでした。ベイズの定理は尤度に事前分布をかけ，周辺尤度で割るという操作によって，事後分布が得られる式，と読むこともできます。\underline{事後分布は\keytermL{確率分布}{かくりつぶんぷ}です}。尤度にある数字をかけて(事後の)確率分布にしていると考えるといいでしょう。
ちなみに事前分布も確率分布ですが，\keytermL{周辺尤度}{しゅうへんゆうど}は確率ではありません。たとえば度数を総度数で割った相対頻度は確率と考えることができますが，それと同じように，分子をとある数字で割って，全体を1.0に整えるための定数なのです。\keytermL{正規化定数}{せいきかていすう}というのはそういう意味ですね。定数倍(正確には定数の逆数)をかけて大きさの調整をしているだけですので，ベイズの式はさらに次のように書き直すことができます。

$$ \mbox{事後分布} \propto \mbox{尤度}\times \mbox{事前分布} $$

ここで$\propto$は比例する，という関係を表しています。最終的には事後分布の形が知りたいのですが，その形を決めているのは分子の尤度と事前分布だけだ，ということになります。

ここで事前分布が\keyterm{一様分布}{いちようぶんぷ}{uniform distribution}であれば，事後分布の形状には影響せず，事後分布は尤度そのものの形を反映することになります。従来のデータ駆動型分析では，データに対して事前の仮定を一切置かないということでしたが，ベイズ流の分析でもそのことは同様に表現できるわけです。

以上がベイズの定理についての簡単な復習でした。しかしベイズの定理自体は，1740年代には明らかになっていたことであり，それがなぜ21世紀の今になって見直されてきたのでしょうか。それには色々な理由がありますが，その1つは最近まで「ベイズ流の分析は絵に描いた餅」だったからです。すなわち理屈ではこのような形になることは明らかだったのですが，実際に計算するのは難しいことが多々ありました。その問題を解決したのが\keyterm{マルコフ連鎖モンテカルロ法}{まるこふれんさもんてかるろほう}{Markov Chain Monte-Carlo method; MCMC}と呼ばれる方法です。


\end{document}